% ═══════════════════════════════════════════════════════════════════════
\section{Systèmes itératifs de fonctions}
% ═══════════════════════════════════════════════════════════════════════

\begin{frame}{La question à l'envers}
  \small Jusqu'ici, on a reconnu des copies à l'intérieur d'une figure
  donnée. Inversons : des contractions $w_1,\dots,w_N$ du plan étant
  données,
  \begin{itemize}
    \item existe-t-il un ensemble $E$ tel que $E = w_1(E)\cup\cdots\cup w_N(E)$ ?
    \item est-il unique ?
    \item comment le construire ?
  \end{itemize}
  \vspace{4pt}
  Réponse : le théorème de Banach, appliqué à $\big(\Comp{\R^2},\dH\big)$.
\end{frame}

\begin{frame}{Définition 3.3 --- Système itératif de fonctions}
  \begin{Enonce}{Définition 3.3 --- Système itératif de fonctions}
    Un \defi{système itératif de fonctions} (SIF) du plan est une famille
    finie $\{w_1,\dots,w_N\}$ de contractions $w_i:\R^2\to\R^2$, de
    rapports $k_1,\dots,k_N$. Le nombre $k=\max\{k_1,\dots,k_N\}$ est le
    \defi{rapport de contraction} du SIF. Son \defi{opérateur de
    Hutchinson} associe à un compact non vide $B$ la réunion de ses
    images :
    \[
      W(B) = w_1(B)\cup w_2(B)\cup\cdots\cup w_N(B).
    \]
  \end{Enonce}
  \small
  \begin{itemize}
    \item Dans ce cours, les $w_i$ sont toujours affines.
    \item $W$ agit sur des figures : un « photocopieur à plusieurs
      lentilles » (Barnsley). En anglais : iterated function system (IFS).
  \end{itemize}
\end{frame}

\begin{frame}{Remarque 3.4 --- Le rapport de contraction d'une application affine}
  \begin{Rem}{Remarque 3.4}
    \small Pour $w(\vec x)=A\vec x+\vec t$ : le plus grand facteur par
    lequel $A$ allonge un vecteur, $\max\norme{A\vec u}$ sur les vecteurs
    unitaires.
    \begin{itemize}
      \item Similitude : son rapport. Diagonale $\mat{a}{0}{0}{b}$ :
        $\max\{|a|,|b|\}$ (le rapport $\tfrac23$ de l'exemple 2.10).
      \item $A$ symétrique : la plus grande valeur propre en valeur absolue.
      \item Mais pas en général : le cisaillement
        $\mat{\tfrac12}{1}{0}{\tfrac12}$ n'a que la valeur propre
        $\tfrac12$, et pourtant il envoie $(0,1)$ sur
        $\big(1,\tfrac12\big)$, de norme $\tfrac{\sqrt5}{2}>1$ : ce n'est
        pas une contraction.
      \item Cas général : racine carrée de la plus grande valeur propre
        de $A^{\mathsf T}A$.
    \end{itemize}
  \end{Rem}
\end{frame}

\begin{frame}{Théorème 3.1 --- Première hypothèse de Banach}
  \small Il faut vérifier que $W$ envoie $\Comp{\R^2}$ dans lui-même
  (admis : une contraction est continue, l'image continue d'un compact
  est compacte, une réunion finie de compacts est compacte), puis :
  \begin{Enonce}{Théorème 3.1 --- L'opérateur de Hutchinson est une contraction}
    Soit un SIF de rapport $k$. Son opérateur $W$ est une contraction de
    rapport $k$ de $\big(\Comp{\R^2},\dH\big)$ :
    $\dH\big(W(A),W(B)\big)\leq k\,\dH(A,B)$.
  \end{Enonce}
  \begin{AuTableau}
    \footnotesize Avec les $\varepsilon$-voisinages (définition 2.8) :
    si $A\subseteq\Vois{B}{\varepsilon}$, alors
    $w_i(A)\subseteq\Vois{w_i(B)}{k\varepsilon}$ ; réunir sur $i$
    (annexe C des notes).
  \end{AuTableau}
\end{frame}

\begin{frame}{Théorème 3.2 --- L'espace des compacts est complet}
  \begin{Enonce}{Théorème 3.2 --- Complétude de l'espace des compacts}
    L'espace métrique $\big(\Comp{\R^2},\dH\big)$ est complet.
  \end{Enonce}
  \small (Admis --- dû à Hausdorff ; une démonstration se trouve chez
  Barnsley.)
  \vspace{6pt}

  Les deux hypothèses du théorème de Banach (forme générale,
  remarque 2.11) sont réunies : un espace complet et une contraction.
\end{frame}

\begin{frame}{Théorème 3.3 --- Hutchinson (1981)}
  \begin{Enonce}{Théorème 3.3 --- Hutchinson, 1981}
    \small Soit un SIF du plan de rapport $k$, d'opérateur de Hutchinson $W$.
    \begin{enumerate}
      \item[(i)] Il existe un unique compact non vide $A$ tel que
        $A = W(A) = w_1(A)\cup\cdots\cup w_N(A)$ ;
      \item[(ii)] pour tout compact non vide $B$, la suite
        $B,W(B),W^2(B),\ldots$ converge vers $A$ au sens de $\dH$ ;
      \item[(iii)] $\dH\big(W^n(B),A\big)\leq\dfrac{k^n}{1-k}\,\dH\big(B,W(B)\big)$.
    \end{enumerate}
    L'ensemble $A$ s'appelle l'attracteur du SIF.
  \end{Enonce}
  \begin{AuTableau}
    \small Banach (forme générale) : $W$ contraction (th.~3.1) d'un
    espace complet (th.~3.2).
  \end{AuTableau}
\end{frame}

\begin{frame}{Un point fixe qui est une figure}
  \small
  \begin{itemize}
    \item Le « point fixe » garanti par Banach n'est pas un point du
      plan, mais une figure entière : l'attracteur.
    \item La suite qui y converge n'est pas une suite de points, mais une
      suite de figures.
    \item Le point (ii) dit de plus que la limite ne dépend pas de la
      figure de départ.
    \item Une démonstration courte, parce qu'elle s'appuie sur tout ce
      qui a été construit aux chapitres 1 et 2.
  \end{itemize}
\end{frame}

\begin{frame}{Figure 3.9 --- La figure de départ n'a pas d'importance}
  \FigBox[0.6\textheight]{%
    \begin{tikzpicture}[scale=1.95, fractale/.style={fill=steelblue4!75, draw=sectioncolor, very thin, line join=round}]
      \foreach \dep/\y in {carre/0, motif/-1.3, disque/-2.6}{
        \foreach \n/\x in {0/0, 1/1.25, 2/2.5, 3/3.75, 5/5}{
          \begin{scope}[shift={(\x,\y)}]
            \draw[gray!45, very thin] (0,0) -- (1,0) -- (0.5,0.8660254) -- cycle;
            \input{figures/fractales/hutchinson-\dep-\n.tikz}
          \end{scope}}}
      \foreach \n/\x in {0/0, 1/1.25, 2/2.5, 3/3.75, 5/5}
        \node[etiq, sectioncolor] at (\x+0.5,1.2) {$W^{\n}(B)$};
    \end{tikzpicture}}
  \small Le SIF de Sierpiński appliqué à un carré, au motif « F » et à un
  disque : chaque fois le même attracteur (point (ii)).
\end{frame}

\begin{frame}{Remarque 3.5 --- Les constructions des sections 2 et 3 sont justifiées}
  \begin{Rem}{Remarque 3.5}
    \small
    \begin{itemize}
      \item Le sens de la limite : le SIF de l'exemple 3.4 appliqué à
        $K_0$ donne exactement $K_1$, puis $K_2=W(K_1)$, etc. La courbe de
        von Koch est la limite de $K_n=W^n(K_0)$ au sens de $\dH$. De même
        $S_n=W^n(S_0)$.
      \item L'autosimilarité : un ensemble autosimilaire
        $E=w_1(E)\cup\cdots\cup w_N(E)$ est un point fixe de $W$ ; par
        unicité, c'est l'attracteur.
    \end{itemize}
    Les ensembles autosimilaires sont exactement les attracteurs des SIF
    de similitudes.
  \end{Rem}
\end{frame}

\begin{frame}{Exemple 3.6 --- Combien d'étapes pour le dessin ?}
  \small Sierpiński dans un triangle de $10$~cm de côté, à un dixième de
  millimètre près, soit $\num{0.001}$ fois le côté.
  \begin{itemize}
    \item $S_1\subseteq S_0$ ; le point de $S_0$ le plus loin de $S_1$
      est le centre du triangle retiré :
      $\dH(S_0,S_1) = \tfrac12\cdot\tfrac{\sqrt3}{6} = \tfrac{\sqrt3}{12}$.
    \item Avec $k=\tfrac12$, le point (iii) donne
      \[
        \dH(S_n,S)\leq\frac{(1/2)^n}{1-1/2}\cdot\frac{\sqrt3}{12}
        = \left(\frac12\right)^{\!n-1}\frac{\sqrt3}{12}.
      \]
    \item Environ $\num{0.00113}$ pour $n=8$ et $\num{0.00056}$ pour
      $n=9$ : neuf étapes, soit $3^9=\num{19683}$ triangles.
  \end{itemize}
\end{frame}

\begin{frame}{Remarque 3.6 --- Les points fixes des $w_i$ sont dans l'attracteur}
  \begin{Rem}{Remarque 3.6}
    \small Chaque $w_i$ a un unique point fixe $\overline{x}_i$ (Banach).
    Partant de $a\in A$, la suite $a,w_i(a),w_i^2(a),\ldots$ reste dans
    $A$ (car $w_i(A)\subseteq W(A)=A$) et converge vers
    $\overline{x}_i$ ; $A$ étant fermé, $\overline{x}_i\in A$.
  \end{Rem}
  \small
  \begin{itemize}
    \item Pour Sierpiński : les trois sommets $S_1$, $S_2$, $S_3$.
    \item En appliquant des composées des $w_i$ à ces points, on obtient
      d'autres points de $A$ --- le principe de l'algorithme
      déterministe (section 9).
  \end{itemize}
\end{frame}
